% !TeX spellcheck = en_US
\documentclass[french]{yLectureNote}

\title{Mecanique Analytique}
\subtitle{Physique}
\author{Paulhenry Saux}
\date{\today}
\yLanguage{Français}
%lionels.calmels@cemes.fr
\professor{M.Brut}
\usepackage{graphicx}%----pour mettre des images
\usepackage[utf8]{inputenc}%---encodage
\usepackage{geometry}%---pour modifier les tailles et mettre a4paper
%\usepackage{awesomebox}%---pour les boites d'exercices, de pbq et de croquis ---d\'esactiv\'e pour les TP de PC
\usepackage{tikz}%---pour deiffner + d\'ependance de chemfig
% \usepackage{tabularx}%---pour dimensionner automatiquement les tableaux avec variable X
\usepackage{awesomebox}%---Pour les boites info, danger et autres
\usepackage{menukeys}%---Pour deiffner les touches de Calculatrice
\usepackage{fancyhdr}%---pour les en-t\^ete personnalis\'ees
\usepackage{blindtext}%---pour les liens
\usepackage{hyperref}%---pour les liens (\`a mettre en dernier)
\usepackage{caption}%---pour la francisation de la l\'egende table vers Tableau
\usepackage{pifont}
\usepackage{array}%---pour les tableaux
\usepackage{lipsum}
\usepackage{yFlatTable}
\usepackage{multicol}
\usepackage{tabularx}
\usepackage{booktabs}
\newcommand{\Lim}[1]{\lim\limits_{\substack{#1}}\:}
\renewcommand{\vec}{\overrightarrow}
\newcommand{\N}[0]{\mathbb{N}}
\newcommand{\dd}{\mathrm{d}}
\newcommand{\norm}[1]{||\vec{#1}||}
\newcommand{\fo}{\psi(\vec{r},t)}
\newcommand{\foe}{\psi(\vec{r},t)\*}
\newcommand{\HH}{\hat{H}}
\newcommand{\hb}{\hbar}
\newcommand{\lap}{\nabla^2}
\newcommand{\lapcc}{\frac{\partial^2 }{\partial x^2}+\frac{\partial^2 }{\partial y^2}+\frac{\partial^2 }{\partial z^2}}
\newcommand{\E}{\vec{E}(M)}
\newcommand{\B}{\vec{B}(M)}
\newcommand{\V}{V(M)}
\newcommand{\er}{\vec{e_{\rho}}}
\newcommand{\ep}{\vec{e_{\varphi}}}
\newcommand{\pip}{\pi^+}
\newcommand{\pim}{\pi^-}
\newcommand{\mv}{\mathcal{V}}
\DeclareMathOperator{\Div}{Div}
\newcommand{\rot}{\vec{\mathrm{rot}}}
\newcommand{\grad}{\vec{\mathrm{grad}}}
\newcommand{\qa}{q_{\alpha}}
\newcommand{\qdb}{\dot{q_{\beta}}}
\newcommand{\qb}{q_{\beta}}
%\setcounter{chapter}{3}
\begin{document}
\chapter{Mécanique lagrangienne}
\section{Introduction}
On se place dans un système dans un espace en D dimensions de coordonéees $x_1,\dots, x_n$\marginCritical{Cela ne correspond pas nécéssairement aux coordonnées de l'espace. Si on étudie 2 particules, cela correspond à la somme des degrés de libertés des 2 particules.}
\begin{definition}[Contrainte holonome]
     On étudie une contrainte qui s'écrit de la forme $$f(x_1,\dots,x_n,t) = 0$$
    \begin{itemize}
        \item Si $f$ ne dépend pas du temps, c'est une contrainte holon\^ome scléron\^ome
        \item Si il y a une dépendance selon $t$, elle est dite rhéon\^ome
    \end{itemize}
\end{definition}
\begin{proposition}[Degrés de libertés]
    Un système déterminé par $D$ coordonnées indépendantes mais soumis à $k$ contraintes holonômes indépendantes
a $N = D - k$ degrés de liberté.
\end{proposition}
Il faut donc N coordonnées généralisées indépendante
pour décrire entièrement le système. 
\section{Lois de Newton et d'Euler-Lagrange}
Soit un système de D dimensions et soumis à k contraintes indépendantes. Le système est soumis aux forces conservatrices $\vec{F_1},\dots, \vec{F_n}$, dont la somme vaut $\vec{F_{ext}}$. On suppose que le mouvement se fait sans frottements
\begin{theorem}[Principe d'Almebert]
    Il y a k forces $\vec{F_{C_i}}$ pour satisfaire les contraintes $f_i(\vec{x})=0$ pour un mouvement sans frottement. Alors 
    $$\vec{F_{C_i}}\propto \vec{\nabla}f_i$$
\end{theorem}
\begin{definition}[Lagrangien]
    $$L(q_{\alpha}, \dot{q_{\alpha}}, t) = U-V = \frac12 m\ddot{x}^2-V$$
\end{definition}
\begin{theorem}[Équation d'Euler-Lagrange]
    $$\frac{\dd }{\dd t}(\frac{\partial L}{\partial \dot{q_{\alpha}}})-\frac{\partial L}{\partial q_{\alpha}} = 0$$
\end{theorem}
\end{document}

